% resume-two-page.tex — layout stress: enough content to span two pages, so a
% page break falls inside the body. Verifies that entries, sections, and lists
% break across pages without page numbers appearing and without overfull boxes.
\documentclass[fontsize=11pt, density=standard]{careerdossier-resume}
\CDossierSetup{
  name     = {Ada Lovelace},
  headline = {Data Scientist},
  email    = {ada.lovelace@example.com},
  phone    = {+1 416 555 0142},
  location = {Toronto, Ontario, Canada},
  linkedin = {linkedin.com/in/ada-lovelace},
  github   = {github.com/ada-lovelace},
}
\begin{document}
\MakeCDossierHeader

\CDossierSection{Summary}
Data scientist with a decade of experience across public-sector analytics,
reporting platforms, and reproducible research pipelines.

\CDossierSection{Experience}
\begin{CDossierEntry}[title = {Principal Data Scientist}, organization = {Digital Services Laboratory}, location = {Ottawa, ON}, dates = {2024--Present}]
  \begin{CDossierItemize}
    \item Led the delivery of reporting services supporting high-volume public programs and internal operational dashboards.
    \item Consolidated several ad hoc analyses into a shared, governed pipeline, reducing duplicated maintenance across program teams.
    \item Established automated validation for data quality, reproducibility, and archival output across every published dataset.
    \item Mentored analysts and defined reusable standards for testing, release management, and technical documentation.
  \end{CDossierItemize}
\end{CDossierEntry}
\begin{CDossierEntry}[title = {Senior Data Scientist}, organization = {Civic Technology Research Group}, location = {Toronto, ON}, dates = {2021--2024}]
  \begin{CDossierItemize}
    \item Designed models and processing pipelines that improved reliability during seasonal demand peaks and large reporting cycles.
    \item Introduced contract tests, deployment checks, and operational dashboards that reduced production regressions.
    \item Led technical discovery with program managers, documented trade-offs, and delivered phased modernization plans.
  \end{CDossierItemize}
\end{CDossierEntry}
\begin{CDossierEntry}[title = {Data Scientist}, organization = {Public Data Analytics Incorporated}, location = {Waterloo, ON}, dates = {2018--2021}]
  \begin{CDossierItemize}
    \item Built data-import services and validation tools for public reporting datasets in CSV, JSON, and XML formats.
    \item Improved processing performance through batching, indexing, and query optimization while preserving auditability.
    \item Created internal libraries for schema validation, error reporting, and automated documentation generation.
  \end{CDossierItemize}
\end{CDossierEntry}

\CDossierSection{Selected Projects}
\begin{CDossierEntry}[title = {Accessible Public Correspondence Platform}, dates = {2025--2026}]
  \begin{CDossierItemize}
    \item Defined a versioned content model supporting conditional sections, reusable clauses, and localization.
    \item Added automated checks for required content, document metadata, reading order, and archival conformance.
  \end{CDossierItemize}
\end{CDossierEntry}
\begin{CDossierEntry}[title = {Open Data Submission Gateway}, dates = {2024--2025}]
  \begin{CDossierItemize}
    \item Implemented schema validation, duplicate detection, detailed error reports, and resumable processing.
    \item Added administrative review tools and audit records to support corrections, approvals, and publication.
  \end{CDossierItemize}
\end{CDossierEntry}
\begin{CDossierEntry}[title = {Operational Reporting Modernization}, dates = {2023--2024}]
  \begin{CDossierItemize}
    \item Replaced spreadsheet-driven monthly reporting with a governed pipeline and reproducible publication process.
    \item Standardized metric definitions and introduced automated reconciliation between source and published summaries.
  \end{CDossierItemize}
\end{CDossierEntry}

\CDossierSection{Education}
\begin{CDossierEntry}[title = {Ph.D. in Computer Science}, organization = {University of Toronto}, location = {Toronto, ON}, dates = {2021--2026}]
\end{CDossierEntry}
\begin{CDossierEntry}[title = {M.Sc. in Software Engineering}, organization = {McGill University}, location = {Montréal, QC}, dates = {2019--2021}]
\end{CDossierEntry}
\begin{CDossierEntry}[title = {B.A.Sc. in Computer Engineering, Honours}, organization = {University of Waterloo}, location = {Waterloo, ON}, dates = {2014--2019}]
\end{CDossierEntry}

\CDossierSection{Certificates}
\begin{CDossierEntry}[title = {Professional Certificate in Accessible Document Production}, organization = {Standards and Accessibility Institute}, dates = {2026}]
\end{CDossierEntry}
\begin{CDossierEntry}[title = {Cloud Architecture and Reliability Engineering}, organization = {Canadian Technology Learning Consortium}, dates = {2025}]
\end{CDossierEntry}

\end{document}
